\documentclass[11pt]{article}
\usepackage[utf8]{inputenc}
\usepackage[spanish, es-tabla]{babel}
\usepackage{amsmath, amssymb}
\usepackage{booktabs}
\usepackage{enumitem}
\usepackage[margin=2.5cm]{geometry}

\setlength{\parindent}{0pt}
\setlength{\parskip}{0.65em}

\title{Econometría ICS-294\\Guía de ejercicios: Clase 4}
\author{Francisco Alfaro Medina}
\date{}

\begin{document}
\maketitle

\section*{Objetivo}

Esta guía busca practicar el sesgo por variables omitidas, su relación con el supuesto $E[\varepsilon\mid X]=0$, la fórmula del sesgo y la interpretación de regresiones con controles.

\section*{Ejercicios}

\begin{enumerate}[label=\textbf{Ejercicio \arabic*.}, leftmargin=*]

\item \textbf{El supuesto clave.}

Considere el modelo:
\[
y_i=\alpha+\beta x_i+\varepsilon_i.
\]

\begin{enumerate}[label=\alph*)]
    \item ¿Qué significa el supuesto $E[\varepsilon_i\mid x_i]=0$?
    \item ¿Por qué este supuesto es necesario para interpretar $\beta$ causalmente?
    \item Dé un ejemplo de una variable que podría quedar dentro de $\varepsilon_i$ y estar correlacionada con $x_i$.
\end{enumerate}

\item \textbf{Consultoría y productividad.}

Una investigadora estima:
\[
productividad_i=\alpha+\beta consultoria_i+\varepsilon_i,
\]
donde $consultoria_i=1$ si la empresa contrató asesoría externa.

\begin{enumerate}[label=\alph*)]
    \item ¿Qué interpretación descriptiva tiene $\hat{\beta}$?
    \item ¿Por qué podría no ser un efecto causal?
    \item Mencione tres variables omitidas que podrían afectar la productividad y también la decisión de contratar consultoría.
    \item Para una de esas variables, explique si esperaría que el sesgo sea positivo o negativo.
\end{enumerate}

\item \textbf{Derivación del sesgo.}

Suponga que el modelo verdadero es:
\[
y=\alpha+\beta_1x_1+\beta_2x_2+u,
\]
con $E[u\mid x_1,x_2]=0$. Además:
\[
x_2=\delta_0+\delta_1x_1+v.
\]

Si se estima el modelo omitido:
\[
y=\tilde{\alpha}+\tilde{\beta}_1x_1+\varepsilon,
\]
demuestre que:
\[
\tilde{\beta}_1=\beta_1+\beta_2\delta_1.
\]

\item \textbf{Cuándo hay sesgo.}

Use la fórmula:
\[
Sesgo=\beta_2\delta_1.
\]

Indique si hay sesgo y su signo en cada caso.

\begin{center}
\begin{tabular}{ccc}
\toprule
Caso & $\beta_2$ & $\delta_1$ \\
\midrule
A & $>0$ & $>0$ \\
B & $>0$ & $<0$ \\
C & $<0$ & $>0$ \\
D & $<0$ & $<0$ \\
E & $=0$ & $>0$ \\
F & $>0$ & $=0$ \\
\bottomrule
\end{tabular}
\end{center}

\item \textbf{Educación, salario e inteligencia.}

Suponga que interesa estimar el efecto causal de la educación sobre el salario. La inteligencia no se observa y se omite del modelo.

\begin{enumerate}[label=\alph*)]
    \item Defina $x_1$, $x_2$ e $y$ en este ejemplo.
    \item ¿Qué signo esperaría para $\beta_2$, el efecto de inteligencia sobre salario?
    \item ¿Qué signo esperaría para $\delta_1$, la relación entre inteligencia y educación?
    \item ¿El coeficiente de educación estaría sobreestimado o subestimado? Explique.
\end{enumerate}

\item \textbf{Cálculo del sesgo.}

Suponga que el efecto verdadero de $x_1$ sobre $y$ es $\beta_1=2$. La variable omitida $x_2$ afecta a $y$ con $\beta_2=3$, y la relación entre $x_2$ y $x_1$ es $\delta_1=0{,}5$.

\begin{enumerate}[label=\alph*)]
    \item Calcule el sesgo por variable omitida.
    \item Calcule $\tilde{\beta}_1$.
    \item ¿El modelo omitido sobreestima o subestima el efecto verdadero?
\end{enumerate}

\item \textbf{Cambio al agregar controles.}

Se estiman dos modelos para explicar salarios:
\[
\widehat{salario}_i=350000+120000\,educacion_i
\]
y luego:
\[
\widehat{salario}_i=280000+85000\,educacion_i+45000\,habilidad_i.
\]

\begin{enumerate}[label=\alph*)]
    \item ¿Cómo cambia el coeficiente de educación al agregar habilidad?
    \item Si habilidad afecta positivamente al salario, ¿qué sugiere este cambio sobre la relación entre habilidad y educación?
    \item ¿La segunda regresión garantiza una interpretación causal? Justifique.
\end{enumerate}

\item \textbf{Cigarrillos y mortalidad.}

En la clase se discute un ejemplo donde una regresión simple entre consumo de cigarrillos y mortalidad entrega un coeficiente negativo. Luego, al controlar por edad, el coeficiente del consumo de cigarrillos aumenta.

\begin{enumerate}[label=\alph*)]
    \item ¿Por qué podría aparecer un coeficiente negativo en la regresión simple?
    \item Si la edad aumenta la mortalidad, ¿qué signo debe tener la relación entre edad y consumo de cigarrillos para producir ese sesgo?
    \item Explique la intuición económica o demográfica detrás de ese signo.
\end{enumerate}

\item \textbf{Controles e interpretación ceteris paribus.}

Explique la diferencia entre estas dos interpretaciones:
\begin{enumerate}[label=\alph*)]
    \item ``Las personas con más educación ganan más''.
    \item ``Comparando personas con igual habilidad observada, las personas con más educación ganan más''.
\end{enumerate}

¿Cuál corresponde a una regresión con controles? ¿Por qué?

\item \textbf{Variables irrelevantes.}

Suponga que el modelo verdadero es:
\[
y=\alpha+\beta_1x_1+u,
\]
pero se estima:
\[
y=\alpha+\beta_1x_1+\gamma z+u.
\]
donde $z$ no afecta a $y$.

\begin{enumerate}[label=\alph*)]
    \item ¿Incluir $z$ genera sesgo?
    \item ¿Qué efecto puede tener sobre la varianza de los estimadores?
    \item ¿Por qué no siempre conviene agregar todas las variables disponibles?
\end{enumerate}

\item \textbf{Aleatorización y controles.}

Suponga que $x_i$ es un tratamiento asignado aleatoriamente.

\begin{enumerate}[label=\alph*)]
    \item ¿Por qué la aleatorización ayuda a cumplir $E[\varepsilon_i\mid x_i]=0$?
    \item ¿Es necesario agregar controles para eliminar sesgo de selección?
    \item ¿Por qué, aun así, podría ser útil agregar algunos controles?
\end{enumerate}

\item \textbf{Diseño aplicado.}

Elija una pregunta causal de interés, por ejemplo: efecto de capacitación sobre empleo, efecto de becas sobre rendimiento, efecto de teletrabajo sobre productividad o efecto de actividad física sobre salud.

Para su caso:
\begin{enumerate}[label=\alph*)]
    \item Escriba una regresión simple.
    \item Identifique una variable omitida relevante.
    \item Explique si la omitida afecta a $y$.
    \item Explique si la omitida está relacionada con $x$.
    \item Determine el signo esperado del sesgo.
    \item Proponga un control que ayudaría a reducir el sesgo.
\end{enumerate}

\end{enumerate}

\section*{Preguntas de cierre}

\begin{enumerate}[label=\arabic*.]
    \item ¿Cuáles son las dos condiciones para que una variable omitida genere sesgo?
    \item ¿Por qué el signo del sesgo depende de $\beta_2$ y $\delta_1$?
    \item ¿Qué cambia en la interpretación de un coeficiente cuando agregamos controles?
    \item ¿Por qué agregar más controles no siempre mejora el modelo?
\end{enumerate}

\end{document}
